\documentclass[10pt]{article}

\usepackage[utf8]{inputenc}

\usepackage[T1]{fontenc}

\usepackage{amsmath}

\usepackage{amsfonts}

\usepackage{amssymb}

\usepackage[version=4]{mhchem}

\usepackage{stmaryrd}

\usepackage{bbold}



\begin{document}

а главной -



\[

\sum_{k=1}^{+\infty} c_{k} z^{k}

\]



Область сходимости ряда (3) имеет вид:



\[

D^{\prime}=\{z: 0 \leqslant r<|z|<R \leqslant+\infty\}

\]



а формулы (2) переходят в формулы:



\[

c_{k}=-\frac{1}{2 \pi i} \int_{\gamma} z^{k+1} f(z) d z, \quad k=0, \pm 1, \ldots,

\]



где $\gamma=\{z:|z|=\rho, r<\rho<R\}$. В остальном все свойства те же, что и в случае конечного центра $a$.



Применение Л. р. основано главным образом на т е о р е ме Л о рана (P. Laurent, 1843): любая однозначная аналитич. функция $f(z)$ в кольце $D=\{z: 0 \leqslant r<|z-a|<R \leqslant+\infty\}$ представима в $D$ сходящимся Л. р. (1). В частности, в проколотой окрестности $D=\{z: 0<|z-a|<R\}$ изолированной особой точки $a$ однозначного характера аналитич. функция $f(z)$ представима Л. р., к-рый и служит основным инструментом исследования ее поведения в окрестности изолированной особой точки.



Для голоморфных функций $f(z)$ многих комплексных переменных $z=\left(z_{1}, z_{2}, \ldots, z_{n}\right)$ аналогом теоремы Лор а н а можно считать следуюшее прелложение: всякую функцию $f(z)$, голоморфную в произведении $D$ колец $D_{v}=\left\{z_{v} \subset \mathbb{C}: 0 \leqslant r_{v}<\left|z_{v}-a_{v}\right|<R_{v} \leqslant+\infty\right\}$, можно представить в $D$ в виде сходяшегося кратного ряда Лорана.



\[

f(z)=\sum_{|k|=-\infty}^{+\infty} c_{k}(z-a)^{k}

\]



в к-ром суммирование распространяется на все целочисленные мультииндексы



\[

\begin{gathered}

k=\left(k_{1}, k_{2}, \ldots, k_{n}\right), \quad|k|=k_{1}+k_{2}+\ldots+k_{n}, \\

(z-a)^{k}=\left(z_{1}-a_{1}\right)^{k_{1}}\left(z_{2}-a_{2}\right)^{k_{2}} \ldots\left(z_{n}-a_{n}\right)^{k_{n}}, \\

c_{k}=\frac{1}{(2 \pi i)^{n}} \int_{\gamma} \frac{f(z) d z}{(z-a)^{k+n}},

\end{gathered}

\]



где $\gamma-$ произведение окружностей $\gamma_{v}=\left\{z_{v} \in \mathbb{C}: z_{v}=a_{v}+\rho_{v} e^{i t}\right.$; $\left.r_{v}<\rho_{v}<R_{v}, 0 \leqslant t \leqslant 2 \pi\right\}, v=1,2, \ldots$, . Область сходимости ряда (4) логарифмически выпуклая. Применение кратных Л. р. (4) ограничено, поскольку при $n \geqslant 2$ голоморфные функции $f(z)$ не могут иметь изолированных особенностей.



Лит.: [1] М а р куш е в и ч А. И., Теория аналитических функций, 2 изд., т. 1, М., 1967, гл. 4; [2] Ш а б а т Б. В., Введение в комплексный анализ, 2 изд., М., 1976, ч. 1, гл. 2, ч. 2, гл. 1. Е. Д. Соломенцев. ЛОРЕНЧА ГАЗ - динамическая система, порожденная движением точечной частицы в пространстве $\mathbb{R}^{d}, d \geqslant 2$, из которого удалено бесконечное число попарно непересекающихся шаров. Частица движется прямолинейно и равномерно до столкновения с шарами и отражается от них по закону упругого удара. Шары (рассеиватели) неподвижны и могут быть расположены как периодически, так и случайно.



Эта модель введена Г. Лоренцом (H. Lorentz, 1905) для описания динамики электронного газа в металлах. Естественная гипотеза состоит в том, что корреляционная функция скоростей для Л. г. убывает как const $t^{-(d / 2+1)}(t-$ время). Л. г. является эргодич. биллиардной системой.



Н. И. Чернов.



ЛОРЕНЦА ГРУППА - группа вещественных линейных однородных преобразований 4-векторов



\[

x=x^{\mu}=\left\{x^{0}, x^{1}, x^{2}, x^{3}\right\}

\]



пространства Минковского $M_{4}$, сохраняющих (индефинитное) скалярное произведение



\[

x y=x^{0} y^{0}-x^{1} y^{1}-x^{2} y^{2}-x^{3} y^{3}=g_{\mu v} x^{\mu} y^{\nu}=x^{\mu} y_{\mu},

\]



где $g=g_{\mu v}-$ метрический тензор в $M_{4}$ (подразумевается суммирование по повторяющимся индексам). Названа по имени Г. Лоренца (H. Lorentz). Являясь подгруппой Пуанкаре группы (группы симметрии пространства-времени в отсутствие гравитации), Л. г. играет фундаментальную роль в релятивистской теории. Инвариантность действия относительно преобразований Л. г. отражает изотропность пространства-времени и влечет за собой сохранение 4-тензора момента (см. Нетер теорема).



Преобразование $\Lambda$ из Л. г. задается вещественной четырехрядной матрицей $\Lambda=\Lambda_{v}^{\mu}$, так что $x^{\mu} \rightarrow x^{\prime \mu}=\Lambda_{v}^{\mu} x_{v}=\Lambda x$. Равенство $\Lambda x \Lambda y=x y$ эквивалентно $\Lambda^{\top} g \Lambda=g\left(\Lambda^{\top}\right.$ транспонирована к $\Lambda$ ) и дает $\operatorname{det} \Lambda= \pm 1,\left|\Lambda_{0}^{0}\right| \geqslant 1$. Л. г. $L$ разбивается на 4 компоненты, связанные между собой в соответствии со знаками $\operatorname{det} \Lambda$ и $\Lambda_{0}^{0}$ :



\[

L=L_{+}^{\uparrow}+L_{+}^{\downarrow}+L_{-}^{\uparrow}+L_{-}^{\downarrow}=L_{+}^{\uparrow}+P T L_{+}^{\uparrow}+P L_{+}^{\uparrow}+T L_{+}^{\uparrow}

\]



Здесь нижний индекс - знак $\operatorname{det} \Lambda$, стрелка $\uparrow(\downarrow)$ отвечает знаку $+(-) \Lambda_{0}^{0} ; P-$ инверсия (отражение) пространства: $(P x)^{0}=x^{0},(P x)^{j}=-x^{j} ; T-$ инверсия времени: $(T x)^{0}=-x^{0}$, $(T x)^{j}=x^{j}, j=1,2,3$. Преобразования с $\operatorname{det} \Delta=1$ называются собственными, с $\Delta_{0}^{0}>0-$ ортохронными. Собственная ортохронная группа $L_{+}^{\uparrow}$ является подгруппой Л. г.



Л. г. - шестипараметрич. группа Ли; в $L_{+}^{\uparrow}$ имеются 3 независимых пространственных вращения $R_{i j}(\alpha)$ на угол $\alpha$ в плоскости $\left(x^{i}, x^{j}\right)$ :



\[

\begin{gathered}

x^{\mu} \rightarrow x^{\mu}, x^{\prime 0}=x^{0}, \\

x^{\prime}=x^{i} \cos \alpha+x^{j} \sin \alpha, \\

x^{\prime}{ }^{j}=x^{j} \cos \alpha-x^{i} \sin \alpha

\end{gathered}

\]



и 3 независимых (частных) Лоренца преобразования - гиперболич. повороты (бусты) $B_{0 k}(\beta)$ на угол $\beta$ в плоскости $\left(x^{0}, x^{k}\right)$



\[

\begin{gathered}

x^{\mu} \rightarrow x^{\prime \mu}, x^{\prime 0}=x^{0} \operatorname{ch} \beta+x^{k} \operatorname{sh} \beta, \\

x^{\prime i}=x^{i}, x^{i j}=x^{j}, x^{\prime k}=x^{k} \operatorname{ch} \beta+x^{0} \operatorname{sh} \beta

\end{gathered}

\]



(здесь $i, j, k=1,2,3$ и их циклич. перестановки). Любой элемент $\Lambda$ из $L_{+}^{\uparrow}$ можно однозначно представить в виде $\Lambda=R B$, где $R$ - пространственное вращение вокруг нек-рой оси, а $B$ - гиперболич. поворот в плоскости $\left(x^{0}, n\right)$, где $\boldsymbol{n}-$ нек-рое направление.



В приложениях важно соответствие между $L_{+}^{\uparrow}$ и группой $S L(2, \mathbb{C})$ комплексных $(2 \times 2)$-матриц с единичным определителем. Каждому $x^{\mu}$ из $M_{4}$ ставится в соответствие эрмитова матрица



\[

X=x^{\mu} \sigma_{\mu}=\left\|\begin{array}{ll}

x^{0}+x^{3} & x^{1}-i x^{2} \\

x^{1}+i x^{2} & x^{0}-x^{3}

\end{array}\right\|,

\]



где $\sigma_{0}$ - единичная $(2 \times 2)$-матрица, $\sigma_{j}$ - Паули матрицы; при этом $x^{\mu} x_{\mu}=\operatorname{det} X$ и $x^{\mu}=(1 / 2) \operatorname{Tr}\left(\sigma_{\mu} X\right)$. Тогда каждому преобразованию $X \rightarrow X^{\prime}=S X S^{+}$, где $S \in S L(2, \mathbb{C})$, отвечает преобразование $\Lambda(S) \in L_{+}^{\uparrow}$, причем



\[

\Lambda_{v}^{\mu}=(1 / 2) \operatorname{Tr}\left(\sigma_{\mu} S \sigma_{v} S^{+}\right)

\]



Это соответствие двузначно: $\Lambda(-S)=\Lambda(S)$; вращениям $R$ отвечают унитарные матрицы $V=S\left(S S^{+}\right)^{-1 / 2}$, бустам $B-$ положительно (либо отрицательно) определенные эрмитовы матрицы $H=\left(S S^{+}\right)^{-1 / 2}$, а разложению $\Lambda=R B-$ разложение $S=V H$. Группа $S L(2, \mathbb{C})$ является универсальной накрывающей Л. г., являясь минимальной односвязной группой, гомоморфной Л. г.



Параметризации Л. г. с помощью углов поворотов отвечает матричное представление ее генераторов $M_{i j}=i R_{i j}^{\prime}(0)$, $M_{0 k}=i B_{0 k}^{\prime}(0)$ (штрих означает здесь производную по углу).





\end{document}